% Compile with LuaLaTeX:
%   lualatex PAPER.tex
% (Japanese support via luatexja / ltjsarticle)
\documentclass[11pt]{ltjsarticle}

\usepackage{luatexja-fontspec}
\usepackage{amsmath,amssymb}
\usepackage{booktabs}
\usepackage{array}
\usepackage{longtable}
\usepackage{graphicx}
\usepackage[margin=25mm]{geometry}
\usepackage{fancyvrb}
\usepackage{enumitem}
\usepackage[hidelinks]{hyperref}

\newcommand{\Ct}{\ensuremath{C_t}}
\newcommand{\norm}[1]{\ensuremath{\lVert #1 \rVert}}

\title{共通目標下における判断基準（V）の不一致は合意品質を低下させるか\\[2pt]
\large ―― HIVC-Dフレームワークのマルチエージェント・シミュレーションによる検証}
\author{HIVC-D検証プロジェクト}
\date{2026年5月}

\begin{document}
\maketitle

\begin{abstract}
対立解消フレームワーク「HIVC-D」は、複数の判断者が共通目標を持っていても、判断基準ベクトル \textbf{V}（価値の重み付け）が不一致であれば合意は真の最適解 \Ct{} から乖離し、結果として成果が低下すると主張する。本研究はこの中核仮説を、二体の探索者が餌場（報酬）を共同探索するマルチエージェント・シミュレーションによって定量的に検証した。V不一致度を3水準（低・中・高）、合意機構を2水準（単純平均=baseline／HIVC-D介入フロー=I共有$\to$V整合チェック$\to$必要時V交渉）に操作し、各セル100試行・計600試行を実施した。

初期実装では仮説H1〜H3はいずれも非有意であったが、内的妥当性に複数の交絡（両エージェントの同一軌跡移動、baseline/HIVC-D間の情報量非対称、\Ct{}定義の不整合）が存在することを発見した。これらを除去する再設計（単一チーム位置モデル、観測平均による情報量対称化、選択肢集合相対の\Ct{}定義）を施した結果、\textbf{測定の信頼性は向上したものの仮説H1〜H3は依然として非有意}であった。因果連鎖を分解すると、理論が想定する第一リンク「V不一致$\to$\Ct{}乖離」が構造的に成立しないことが判明した（線形回帰 slope$=-0.033$, $p=0.49$, $r^2=0.0016$）。一方、仮説H4「I共有後の順位相関 $\rho$ はV不一致度に反比例する」は一貫して強く成立した（$r^2=0.733$, $p<0.001$）が、これはスコア計算のみで閉じる準・自明な関係であった。

結論として、\textbf{「全選択肢が共通目標と整合し、合意が対称な平均で、目標が共有される」レジームにおいては、V不一致は平均操作によって自己相殺され、合意品質を低下させない}。HIVC-D介入が優位性を示すには、Vが系統的に劣選択へバイアスをかける構造（誤誘導情報・非対称合意・目標と直交するV）が必要であると示唆される。
\end{abstract}

\noindent\textbf{キーワード}: 合意形成, 価値の不一致, マルチエージェント・シミュレーション, 集団意思決定, 帰無結果

\section{はじめに (Introduction)}

\subsection{背景とアイデアの起点}

組織やチームにおける意思決定では、「全員が同じゴールを目指しているのに、なぜ合意がうまくいかないのか」という問いがしばしば現れる。HIVC-Dフレームワークはこの現象を、以下の要素に分解して説明する。

\begin{itemize}[nosep]
  \item \textbf{共通目標 G を持つこと}は合意成立の十分条件ではない。
  \item 各判断者が用いる\textbf{判断基準 V（何をどれだけ重視するか）}が異なれば、同じ情報・同じ目標からでも異なる結論に至る。
  \item さらに各判断者が参照する\textbf{環境情報 I}も異なりうる。
\end{itemize}

本研究の出発点となるアイデアは次の一点に集約される：

\begin{quote}
\textbf{共通目標Gの存在にもかかわらず、判断基準Vの不一致それ自体が、集団の合意を真の最適解から乖離させ、成果を毀損する独立の要因となるのではないか。}
\end{quote}

もしこれが正しければ、対立解消の介入は「目標のすり合わせ」よりも「\textbf{判断基準の診断と整合（V整合）}」を優先すべきだという実践的含意が導かれる。HIVC-Dはこの介入順序を「I（情報）の共有 $\to$ V（基準）の整合 $\to$ A（能力）の確認」という診断フローとして規定する。

\subsection{本研究の目的}

本研究は、抽象的な理論主張を反証可能な形に操作化し、シミュレーションによって以下の4仮説を検証することを目的とする。

\begin{longtable}{@{}l p{0.78\linewidth}@{}}
\toprule
仮説 & 内容 \\
\midrule
\textbf{H1} & V不一致度が高いほど、合意の成果（累積報酬）は低下する。 \\
\textbf{H2} & HIVC-D介入フロー（条件B）は単純平均合意（条件A）より高い成果を生む。 \\
\textbf{H3} & V不一致が大きい条件ほど、介入の改善余地（A-B差）が大きい。 \\
\textbf{H4} & 共通情報Iの共有はVに依存せず順位相関 $\rho$ の収束を促す。 \\
\bottomrule
\end{longtable}

\section{HIVC-Dフレームワーク (The HIVC-D Framework)}

本節は、HIVC-Dを初めて読む研究者のために、フレームワークの目的・構成要素・主張を自己完結的に説明する。すでに親しんでいる読者は 2.7 の検証スコープへ進んでよい。

\subsection{フレームワークが扱う問題}

HIVC-Dは、\textbf{複数の判断者が「合意」に至るプロセスを分解・診断するための記述的フレームワーク}である。動機となる経験的観察は次の一点に集約される。

\begin{quote}
チームの全員が同じ目標を掲げているにもかかわらず、しばしば合意は形成されず、形成されても誤った結論に至る。
\end{quote}

日常的にはこの失敗が「人間関係の問題」や「コミュニケーション不足」と一括りにされがちである。HIVC-Dの主張は、こうした合意失敗が実は\textbf{少数の識別可能な要因（情報・判断基準・能力）のいずれかのズレ}に還元でき、要因を特定すれば対応する介入を選べる、というものである。すなわちフレームワークの目標は「対立を精神論でなく\textbf{診断可能な構造}として扱う」ことにある。

\subsection{名称と基本要素}

フレームワーク名 \textbf{HIVC-D} は、中核となる構成要素の頭文字（判断者 \textbf{H}、情報 \textbf{I}、判断基準 \textbf{V}、選択肢 \textbf{C}）に、反復性を表す拡張 \textbf{D（Dynamic）} を付したものである。判断状況は以下の要素で記述される。

\begin{Verbatim}[frame=single,fontsize=\small]
E   : 環境(Environment)（外部制約・リソース上限を含む）
H   : 判断者集合(Human) {h1, h2, ...}
G   : 共通目標(Goal)（合意済みであることを前提とする）
P   : 解くべき命題(Problem)（何を決めるのか）
C   : 選択肢集合(Choices)（取りうる手）
C_t : Gに対する真の最適選択肢(Choice True)（存在は仮定するが一意とは限らない）
I   : 環境情報(Information)（各判断者が参照するデータ・事実）
V   : 判断基準(Value)（各判断者が選択肢を評価する際の価値の重み付け）
A   : 各hの推定能力(Ability)（情報処理・判断の精度上限）
\end{Verbatim}

直感的には、各判断者 h は自分の持つ情報 \textbf{I} を、自分の価値観 \textbf{V} で重み付けして選択肢 \textbf{C} を評価し、最良と思う手を選ぶ。全員が真の最適解 \Ct{} を選べば合意は自明だが、現実には I・V・A のいずれかが人によって異なるため、評価が食い違い対立が生じる。

\textbf{重要な前提}: HIVC-Dは共通目標 G が\textit{すでに合意されている}状況を扱う。つまり「ゴールが違うから揉める」という自明なケースは対象外であり、\textbf{「同じゴールなのになぜ揉めるのか」}に焦点を絞る。これが本フレームワークの問題設定を鋭くしている点である。

\subsection{オリジナル版の構造と弱点}

初期のHIVC-Dは「全判断者が理想的に推論すれば必ず\Ct{}に収束し合意する」という静的モデルであった。しかしこの定式化には以下の理論的弱点がある。本研究はこれらを修正した改定版（Dynamic）を検証対象とする。

\begin{longtable}{@{}p{0.30\linewidth} p{0.62\linewidth}@{}}
\toprule
問題 & 具体的欠陥 \\
\midrule
VとIの独立仮定 & 情報がVを変え、VがIの収集方向を変えるという相互依存を無視 \\
トートロジー仮定 & 「全員が理想的なら合意する」は定義であって予測力を持たない \\
静的モデル & 一回限りの判断として扱い、実務の反復性（フィードバック）を無視 \\
能力項の粗さ & 「能力不足」を対策不能なものとして切り捨て \\
合意$\neq$実行 & 合意後のコミットメント・実行が未定義 \\
\bottomrule
\end{longtable}

\subsection{改定版 HIVC-D (Dynamic) の修正仮定}

\begin{longtable}{@{}p{0.34\linewidth} p{0.58\linewidth}@{}}
\toprule
旧仮定 & 修正 \\
\midrule
\Ct{}は一意に存在 & \Ct{}は\textbf{Pareto最適集合}として存在。一意性は保証しない \\
IとVは独立 & $I(t)$ と $V(t)$ は\textbf{相互依存}。IがVを更新し、VがIの収集方向を決める \\
全員理想$\to$合意（トートロジー） & \textbf{削除}し、合意条件を明示的に定義（2.5）に置き換える \\
能力は固定 & Aは訓練・委任・分業によって\textbf{操作可能なパラメータ} \\
\bottomrule
\end{longtable}

\subsection{合意条件の明示定義（トートロジーの置き換え）}

改定版は「理想なら合意」という定義を捨て、合意成立を3つの検証可能な条件で定義する。

\begin{Verbatim}[frame=single,fontsize=\small]
任意の判断者ペア hi, hj ∈ H について以下が成立するとき合意が得られる：

  1. I一致条件 : hi と hj が参照する情報 I の差分が許容誤差 ε 以下
  2. V整合条件 : Vi と Vj の優先順位が一致、または双方が許容できる V* に収束
  3. 能力条件  : 各 h の能力 A が C_t を評価するのに十分
\end{Verbatim}

この3条件は、後述する診断（2.6）の出発点でもある。合意が成立しないとき、3条件のうちどれが破れているかを問うことで原因を切り分けられる。

\subsection{動的モデルと診断・介入}

改定版の本質は、判断を\textbf{一回限りでなく環境フィードバックを伴う反復プロセス}として捉える点にある。

\begin{Verbatim}[frame=single,fontsize=\small]
t=0: 共通目標 G を合意
      |
Step 1: I(t) を収集     -- V が I の収集方向を規定する
      |
Step 2: V(t) を更新     -- 収集した I によって V が修正されうる
      |
Step 3: C(t) を評価     -- 現在の I(t), V(t), A(t) で C_t を推定
      |
Step 4: 合意判定        -- 3条件(I一致・V整合・能力)を満たすか？
      |- YES -> 実行・コミット -> 環境フィードバック E(t+1), I(t+1) -> 次サイクルへ
      `- NO  -> 診断 -> 介入（下表）
\end{Verbatim}

合意が成立しない場合、まず\textbf{「同一の $I(t)$ を与えたとき合意するか？」}を問うことで原因を診断する。

\begin{itemize}[nosep]
  \item \textbf{YES（同じ情報なら合意する）} $\to$ 原因は \textbf{V不一致}。
  \item \textbf{NO（同じ情報でも合意しない）} $\to$ 原因は \textbf{I不一致} または \textbf{能力不足}。
\end{itemize}

診断結果に応じて、コストの低い介入から順に適用する。

\begin{longtable}{@{}l l l p{0.42\linewidth}@{}}
\toprule
介入手段 & 対象原因 & コスト & 備考 \\
\midrule
I共有 & I不一致 & 低 & 最初に試みる \\
V交渉 & V不一致 & 中 & Gへの整合を軸に共通 $V^*$ を探す \\
C列挙 & I・V複合 & 中 & 双方のVを部分充足する新たな選択肢を探索 \\
A強化 & 能力不足 & 高 & 訓練・委任・外部専門家の導入 \\
ランダムC & 時間切れ & --- & \Ct{}への収束を\textbf{明示的に諦める}決断 \\
\bottomrule
\end{longtable}

この「\textbf{診断順序 I $\to$ V $\to$ A}」こそがHIVC-Dの実践的中核である。情報のズレ（I）は安価に解消できるので最初に潰し、それでも合意しなければ価値観のズレ（V）を交渉し、最後に能力（A）の問題を扱う。

\subsection{本研究が検証する主張とそのスコープ}

以上のフレームワーク全体のうち、本研究が定量検証するのは次の二点に限定される。

\begin{enumerate}[nosep]
  \item \textbf{中核仮説（V不一致の害）}: 共通目標Gの下でも、Vの不一致は合意を真の最適解\Ct{}から乖離させ、成果を低下させる。
  \item \textbf{診断順序の優位（HIVC-D介入）}: 「I共有 $\to$ V整合チェック $\to$ 必要時V交渉」という介入フローは、単純な妥協（平均）より優れた成果を生む。
\end{enumerate}

能力Aの不均一、多人数化、C列挙やA強化といった介入、および合意後の実行・フィードバックは本研究のスコープ外とし、Aは均一に固定する。

理論は中核仮説を次の\textbf{二段の因果連鎖}として予測する。本研究の方法論的工夫は、この連鎖を分解して各リンクを個別に測定する点にある。

\begin{Verbatim}[frame=single,fontsize=\small]
V不一致度 --(リンク1)--> 合意ベクトルのC_t乖離 --(リンク2)--> 累積報酬の低下
\end{Verbatim}

リンク1が成立しなければ、たとえリンク2（乖離が成果を下げること）が真でも、Vは成果に影響しない。第4節で示すとおり、本研究の主要な発見はこのリンク1の不成立である。

\subsection{フレームワーク要素のシミュレーション上の対応}

抽象的な要素を、餌場探索シミュレーションへ次のように操作化した（詳細は第3節）。

\begin{longtable}{@{}l l p{0.52\linewidth}@{}}
\toprule
記号 & 名称 & シミュレーション上の対応 \\
\midrule
E    & 環境 & フィールド（餌・標識の配置） \\
H    & 判断者集合 & 二人の探索者 \\
G    & 共通目標 & 累積報酬の最大化 \\
C    & 選択肢集合 & 進行方向（連続空間のベクトル） \\
\Ct{} & 真の最適選択肢 & 報酬効率が最大の方向（ノイズなし理想解） \\
I    & 環境情報 & 各探索者が観測した標識の集合 \\
V    & 判断基準 & 重みベクトル $(w_c, w_d, w_p)$ \\
A    & 推定能力 & 本実験では均一に固定（制御変数） \\
\bottomrule
\end{longtable}

判断基準 $V = (w_c, w_d, w_p)$ は、各探索者が標識を評価する際の重みであり、それぞれ「進行方向との一貫性 ($c$)」「距離効率 ($d$)」「報酬量 ($p$)」への選好を表す（合計1に正規化）。合意の3条件は本シミュレーションで次のように操作化される。

\begin{itemize}[nosep]
  \item \textbf{I一致条件}: HIVC-D条件では両探索者の観測を共有（\texttt{shared\_obs}）し強制的に満たす。
  \item \textbf{V整合条件}: 共通情報上のスコア順位のSpearman順位相関 $\rho$ が閾値（\texttt{RHO\_AGREE\_THRESHOLD} $= 0.6$）以上であること。$\rho < 0.6$ のとき「V不整合」と判定しV交渉（$v^* = (v_1 + v_2)/2$、双方が等分に譲歩）に移行する。
  \item \textbf{能力条件}: 本実験では均一に固定しスコープ外とする。
\end{itemize}

\section{実験設計 (Methods)}

\subsection{環境}

$100\times100$ の連続二次元フィールドに、報酬量 $[1.0, 10.0]$ の餌を20個、標識を50個ランダム配置した。各標識は最近傍の餌を指し示し、観測時に方向ノイズ（$\sigma_\theta=0.3$ rad）、距離ノイズ（$\sigma_d=0.2$、乗法的）、報酬ノイズ（$\sigma_p=1.0$、加法的）を独立に付与する。探索者は毎ターン最近傍 $K=5$ 枚の標識を観測し、合意した方向へ \texttt{MOVE\_STEP}$=2.0$ だけ移動する。餌獲得判定半径は2.0。1試行は \texttt{T\_MAX}$=200$ ターン。

\subsection{独立変数}

\textbf{(a) V不一致度（3水準）}

\begin{longtable}{@{}l c c c@{}}
\toprule
条件 & $V_1$ & $V_2$ & $\norm{V_1-V_2}_2$ \\
\midrule
V-Low  & $(0.33, 0.33, 0.33)$ & $(0.40, 0.30, 0.30)$ & 0.082 \\
V-Mid  & $(0.60, 0.20, 0.20)$ & $(0.20, 0.20, 0.60)$ & 0.566 \\
V-High & $(0.80, 0.10, 0.10)$ & $(0.10, 0.10, 0.80)$ & 0.990 \\
\bottomrule
\end{longtable}

V-Mid・V-Highでは $w_d$（距離重み）が両探索者で同一であり、不一致は主に「方向一貫性重視 ($w_c$)」対「報酬量重視 ($w_p$)」の対立として現れる。

\textbf{(b) 合意機構（2水準）}

\begin{itemize}[nosep]
  \item \textbf{条件A（baseline）}: 各探索者が独立に希望ベクトルを計算し、その単純平均を進行方向とする。I共有もV交渉も行わない。
  \item \textbf{条件B（HIVC-D）}: \textcircled{1} I共有（両探索者の観測を統合）$\to$ \textcircled{2} V整合チェック（$\rho$ 算出）$\to$ \textcircled{3} $\rho\ge0.6$ なら統合情報上で希望ベクトルを平均、$\rho<0.6$ ならV交渉（$v^*$ で希望ベクトルを再計算）。
\end{itemize}

$3\times2=6$ セル、各100試行、計600試行。乱数は \texttt{seed = 42 + trial\_id} で再現性を確保した。

\subsection{従属変数}

\begin{itemize}[nosep]
  \item \textbf{累積報酬}（主要成果指標）
  \item \textbf{\Ct{}近似率}（合意ベクトルと\Ct{}の内積、リンク1の直接測定）
  \item \textbf{順位相関 $\rho$}（V整合度、条件Bのみ）
  \item \textbf{平均合意コスト}（V交渉発生率）
\end{itemize}

\subsection{統計手法}

\begin{itemize}[nosep]
  \item H1: baselineのみでV-Low/Mid/Highの一元配置ANOVA、およびV距離を連続変数とする線形トレンド回帰。
  \item H2: 各V条件でWelchのt検定（baseline vs hivc）、3条件の多重比較にBonferroni補正（$\alpha=0.0167$）。
  \item H3: 二元配置ANOVA（V条件$\times$機構）の交互作用項。
  \item H4: $\rho_{\text{before}}$（V交渉前）をV不一致度で線形回帰。
\end{itemize}

\subsection{実装の反復改良と内的妥当性の確保}

初期実装で全仮説が非有意となったため、内的妥当性の体系的レビューを実施し、以下の交絡を特定・除去した。この過程自体が本研究の方法論的貢献である。

\begin{longtable}{@{}p{0.20\linewidth} p{0.40\linewidth} p{0.32\linewidth}@{}}
\toprule
発見した交絡 & 問題 & 修正 \\
\midrule
\textbf{同一軌跡移動} & 両エージェントが毎ターン同一ベクトルで移動し、初期座標差が定数保存。実質的に擬似二体問題化 & 単一チーム位置モデルへ移行。I不一致を観測ノイズとして純化 \\
\textbf{情報量の非対称} & baselineはK枚、HIVC-Dは最大2K枚を使用し、比較が「機構の差」でなく「情報量の差」を測定 & 同一標識をインデックス対応で観測平均し、双方K枚に対称化 \\
\textbf{\Ct{}定義の不整合} & \Ct{}をフィールド全体の最良餌（遠方含む）とし、探索者の可視選択肢と乖離。\Ct{}近似率が系統的に負値（$-0.27$） & 可視標識集合C内のノイズなし最良方向として\Ct{}を再定義 \\
\textbf{スコア距離項のスケール} & 任意定数による正規化 & \texttt{FOOD\_RADIUS}を単位とする物理的正規化 \\
\textbf{フィールドリセット} & \texttt{seed+t}でリセット時点に依存 & \texttt{seed+reset\_count$\times$100003}で決定論化 \\
\textbf{方向の重複判定} & 0.1rad閾値が任意、$\pm\pi$折り返し未対応 & 循環平均に置換 \\
\bottomrule
\end{longtable}

これらの修正後、\Ct{}近似率の符号が \textbf{$-0.27 \to +0.53$} に正常化し、因果連鎖を信頼できる形で分解測定可能となった。

\section{結果 (Results)}

\subsection{主要仮説の検定（最終モデル、$n=100$/セル）}

\begin{longtable}{@{}l l l l@{}}
\toprule
仮説 & 統計量 & $p$値 & 判定 \\
\midrule
H1（ANOVA） & $F=0.196$, $\eta^2=0.001$ & 0.822 & 非有意 \\
H1（線形回帰） & slope$=0.350$, $r^2=0.000$ & 0.736 & 非有意 \\
H2 V-Low & $t=-0.923$ & 0.357 & 非有意 \\
H2 V-Mid & $t=-0.987$ & 0.325 & 非有意 \\
H2 V-High & $t=-0.383$ & 0.702 & 非有意 \\
H3（交互作用） & $F=0.091$ & 0.914 & 非有意 \\
H4（$\rho$ vs V距離） & slope$=-0.661$, $r^2=0.733$ & $<0.001$ & \textbf{有意} \\
\bottomrule
\end{longtable}

\subsection{条件別の累積報酬}

\begin{longtable}{@{}l l c c@{}}
\toprule
V条件 & 機構 & 平均 & 標準偏差 \\
\midrule
V-Low  & baseline & 7.06 & 6.70 \\
V-Low  & hivc     & 6.21 & 6.30 \\
V-Mid  & baseline & 6.80 & 6.20 \\
V-Mid  & hivc     & 5.98 & 5.65 \\
V-High & baseline & 7.39 & 7.08 \\
V-High & hivc     & 7.03 & 6.44 \\
\bottomrule
\end{longtable}

V不一致度の増加に伴う系統的な報酬低下（H1）も、HIVC-Dの優位性（H2）も観測されなかった。むしろHIVC-Dは全条件でわずかに低い平均を示した。

\subsection{因果連鎖の分解}

理論の二段連鎖を個別に測定した結果（baseline条件）：

\begin{Verbatim}[frame=single,fontsize=\small]
リンク1: V不一致 -> C_t乖離     slope=-0.033, p=0.49, r^2=0.0016   -- 不成立
リンク2: C_t乖離 -> 累積報酬    slope=-1.58,  p=0.21               -- 弱い
\end{Verbatim}

\Ct{}近似率は全V条件でほぼ一定であった（V-Low: 0.558, V-Mid: 0.530, V-High: 0.528）。すなわち、\textbf{V不一致度を約12倍（$0.082\to0.990$）に増大させても、合意ベクトルの\Ct{}からの乖離は変化しなかった}。

\subsection{介入効果の直接測定}

HIVC-Dが\Ct{}近似率を改善するかを直接検定した結果、いずれのV条件でも有意な改善は見られなかった（V-Low: $\Delta=-0.001$ $p=0.98$、V-Mid: $\Delta=+0.022$ $p=0.62$、V-High: $\Delta=-0.012$ $p=0.78$）。

\subsection{ハイパーパラメータのグリッドサーチ}

\texttt{RHO\_AGREE\_THRESHOLD}（4水準）$\times$ \texttt{SIGMA\_THETA}（3水準）$\times$ \texttt{K\_SIGNS}（3水準）の36設定$\times$各30試行を探索した結果、H1有意は0/36、H2が1条件でも有意となるのは1/36であった。観測ノイズ $\sigma_\theta$ が最小（0.1）かつ標識数 \texttt{K\_SIGNS} が最大（10）のときにHIVC-Dの優位性が最大化する傾向はあったが、いずれも統計的に有意な水準には達しなかった。

\section{考察 (Discussion)}

\subsection{なぜV不一致は無害だったのか}

最も重要な知見は、理論の第一リンク「V不一致$\to$\Ct{}乖離」が本実験のレジームで構造的に成立しないことである。その機序は以下の三条件の組み合わせに帰着する。

\begin{enumerate}[nosep]
  \item \textbf{全選択肢が共通目標と整合}: すべての標識は実在の餌を（ノイズ付きで）指し示しており、誤誘導する選択肢が存在しない。
  \item \textbf{対称な平均合意}: 合意は二つの希望ベクトルの単純平均である。
  \item \textbf{目標の共有}: 両探索者は同一の累積報酬最大化を目指す。
\end{enumerate}

この条件下では、異なるVで重み付けされた二つの希望ベクトルは、それぞれ異なる方向に偏るものの、\textbf{平均操作によって偏りが相殺}され、結果は不偏な平均に収束する。たとえばV-High条件では一方が「方向一貫性」、他方が「報酬量」を重視するが、両者を平均すると両選好を程よく満たす中庸な方向となり、\Ct{}への近似度は単一基準の場合と変わらない。すなわち\textbf{解くべき対立が存在しない}ため、HIVC-D介入も発揮すべき効果を持たない。

\subsection{H4の有意性の解釈}

H4（$\rho$ vs V距離）のみが一貫して強く成立したが、これは理論の検証というより\textbf{準・自明な関係}である。$\rho$ は共通情報上のスコア順位の一致度であり、エージェントの移動や報酬を一切経由せずスコア計算だけで決定される。Vが異なれば順位が変わるのは定義上ほぼ必然であり、この関係が成立しても「Vが\Ct{}乖離や成果に影響する」という主張は支持されない。H4を中核証拠として扱うことには注意を要する。

\subsection{帰無結果の価値と方法論的含意}

本研究は仮説を支持しない結果に終わったが、これは\textbf{信頼できる帰無結果}である。反復的な内的妥当性レビューによって交絡を除去した上での結論であり、「仮説を有意化するための恣意的なパラメータ調整（p-hacking）」を行わなかった。むしろ交絡の除去過程で、初期実装の見かけ上の挙動が複数のアーティファクトに起因していたことを明らかにした点に方法論的価値がある。

\subsection{仮説が成立しうるレジーム}

理論を救済しうる条件は、5.1で挙げた前提のいずれかを破ることである。

\begin{enumerate}[nosep]
  \item \textbf{誤誘導/敵対的情報の導入}: 一部の標識が虚偽の餌情報を提示する環境。Vが「報酬量」に強く偏った探索者は虚偽の高報酬標識に誘導され、平均合意でも偏りが残存しうる。
  \item \textbf{非対称合意}: 一方の探索者が支配的な重みを持つ合意。相殺が働かず、支配側のVバイアスが直接結果に反映される。
  \item \textbf{目標と直交するV成分}: 報酬最大化と無関係な価値（例: 特定方向への固執）を持つ探索者。この成分は平均しても相殺されず、系統的な\Ct{}乖離を生む。
\end{enumerate}

これらの設定では、HIVC-Dの「I$\to$V診断順序」が単純平均を上回る余地が理論的に生じる。

\section{限界と今後の課題 (Limitations and Future Work)}

\begin{itemize}[nosep]
  \item \textbf{単一レジームの検証}: 本実験は「協調的・整合的・対称的」という特定のレジームのみを扱った。5.4の代替レジームでの検証が次の課題である。
  \item \textbf{V交渉モデルの単純性}: 等分譲歩 $v^*=(v_1+v_2)/2$ のみを実装した。Nashの交渉解など他の交渉アルゴリズムの効果は未検証である。
  \item \textbf{二者・固定能力}: 判断者は二名、推定能力Aは均一に固定した。多人数化やA不均一の影響はスコープ外である。
  \item \textbf{報酬指標のノイズ}: 累積報酬は経路依存性が高く分散が大きい。\Ct{}近似率のようなより直接的な指標を主要従属変数とする設計が、検出力の観点で望ましい。
\end{itemize}

\section{結論 (Conclusion)}

共通目標下における判断基準Vの不一致が合意品質を低下させるというHIVC-Dの中核仮説を、マルチエージェント・シミュレーションで検証した。内的妥当性を確保した最終モデルにおいて、\textbf{V不一致は合意ベクトルの\Ct{}乖離を引き起こさず（理論の第一リンクが不成立）、累積報酬にも影響しなかった}。この帰無結果は、全選択肢が目標と整合し合意が対称な平均であるレジームでは、V不一致が平均操作によって自己相殺されることに起因する。HIVC-D介入が優位性を発揮するには、Vが系統的に劣選択へバイアスをかける環境構造（誤誘導情報・非対称合意・目標と直交するV）が必要であると結論づけられる。本研究は、理論の適用範囲を限定し、その境界条件を明示した点に貢献がある。

\section*{コード・データの公開 (Code and Data Availability)}

本研究で用いた全シミュレーションコード、統計検定、グリッドサーチ、および生成された結果（図・集計統計量）は、以下のリポジトリで MIT ライセンスのもと公開している。

\begin{quote}
\textbf{GitHub}: \url{https://github.com/neipia271828/hivc-d-verification}
\end{quote}

リポジトリには本論文（\texttt{PAPER.md} / \texttt{PAPER.tex}）、再現手順（\texttt{README.md}）、依存関係（\texttt{requirements.txt}）、単体テスト一式を含む。全乱数はシードで制御されており、結果は完全に再現可能である。

\section*{付録: 再現方法}

リポジトリを取得し、依存関係をインストールした上で実行する。

\begin{Verbatim}[frame=single,fontsize=\small]
git clone https://github.com/neipia271828/hivc-d-verification.git
cd hivc-d-verification
pip install -r requirements.txt
cd hivc_sim

# 単体テスト
python -m pytest tests/ -v

# 本実験（全600試行）
python main.py --trials 100

# ハイパーパラメータ・グリッドサーチ
python grid_search.py --trials 30
\end{Verbatim}

出力: \texttt{results/raw/trials.csv}（試行別生データ）、\texttt{results/summary/stats.csv}（検定統計量）、\texttt{results/figures/}（図）、\texttt{results/grid\_search/}（探索結果）。

\end{document}
